% ========================================================
% ANEXOS
% ========================================================
\section{Video de sustentación}\label{anexo:video}

El enunciado pide un video de hasta seis minutos, publicado en la nube, en el que cada integrante
demuestra conocimiento del tema, presenta los resultados y discute críticamente las limitaciones, la
escalabilidad y las posibles mejoras.

\begin{itemize}
  \item \textbf{URL:} \pendiente{enlace público al video}
  \item \textbf{Duración:} \pendiente{mm:ss, como máximo 6:00}
  \item \textbf{Partes:} \pendiente{quién presenta qué, con el minuto en que empieza cada parte}
\end{itemize}

\section{Informe de GAPs en Markdown}\label{anexo:gaps}

El informe de GAPs de la Sección~\ref{sec:gaps} está en el repositorio, en la rama principal:
\url{https://github.com/R0SEWT/concurrente/blob/main/tp/docs/gaps-ia.md}.

\section{Declaración de uso de inteligencia artificial}\label{anexo:ia}

El sílabo del curso incorpora la ingeniería de prompts como herramienta, siempre contrastada con la
teoría. En este trabajo se usó IA generativa en las tareas que siguen, y todo lo producido con ella
fue revisado por el grupo y validado contra el código, los tests y las salidas de Spin, que son la
fuente de las cifras de este informe.

\pendiente{lista revisada por el grupo: qué herramienta y modelo se usó en cada tarea (código,
modelo Promela, redacción del informe, análisis de GAPs) y qué se verificó a mano. Partir del
registro de uso de IA (bead r0p.1)}
